\section{Concurrency Theory}

\subsection{Linearizability}

\begin{definition}
A \textbf{history} is a partially ordered set of invocations and responses of operations in a distributed program.
\end{definition}

\begin{definition}
A history $H$ is \textbf{linearizable} if it can be extended to a history $G$ by
\begin{itemize}
    \item appending zero or more responses to pending invocations that took effect
    \item discarding zero or more pending invocations that did not take effect
\end{itemize}
such that $G$ is equivalent to a legal sequential history $S$ with $\to_G \subseteq \to_S$.
\end{definition}

\begin{theorem}
History $H$ is linearizable $\iff$ for every object $x$, $H | x$ is linearizable.
\end{theorem}

\begin{remark}
Linearizability is a \textbf{local property}: it can be proven in isolation for each object and then composed.
\end{remark}

\subsection{Sequential Consistency}

\begin{definition}
A history $H$ is \textbf{sequentially consistent} if it can be extended to a history $G$ by
\begin{itemize}
    \item appending zero or more responses to pending invocations that took effect
    \item discarding zero or more pending invocations that did not take effect
\end{itemize}
such that $G$ is equivalent to a legal sequential history $S$ (no real-time order constraint, only program order of each thread).
\end{definition}

\begin{theorem}
Sequential consistency is \textbf{not a local property}, thus we lose composability.
\end{theorem}

\begin{remark}
\textbf{Quiescence consistency} requires non-overlapping methods to take effect in real-time order (periods of quiescence separate operations).
\end{remark}

\subsection{Consensus}

\begin{definition}
A \textbf{consensus object} has a function $decide(value)$ that returns the same value for all threads, where the value is one of the values $v_i$ proposed by the threads.
\end{definition}

\begin{definition}
The \textbf{Consensus Number} is the maximum number of nodes that can solve consensus validly, sequentially consistently and wait-free.
\end{definition}

\begin{theorem}
\textbf{Consensus Numbers:}
\begin{itemize}
    \item Atomic Registers: consensus number $= 1$
    \item Test-and-Set, Get-and-Set: consensus number $= 2$
    \item Compare-and-Swap (CAS): consensus number $= \infty$
\end{itemize}
\end{theorem}

\begin{remark}
\textbf{CAS has infinite consensus number} because we can implement a valid, sequentially consistent, and wait-free consensus object using CAS:
\end{remark}

\begin{codeexample}
class CASConsensus {
    private final int FIRST = -1;
    private AtomicInteger r = AtomicInteger(FIRST);
    private AtomicIntegerArray proposed;

    public int decide(int value) {
        int i = ThreadID.get();
        proposed.set(i, value);
        if (r.compareAndSet(FIRST, i)) {
            return proposed.get(i);
        } else {
            return proposed.get(r.get());
        }
    }
}
\end{codeexample}

\begin{remark}
There is \textbf{no wait-free FIFO Queue using atomic registers}: a wait-free FIFO Queue can solve 2-consensus, contradicting the consensus number of 1 for atomic registers.
\end{remark}

\subsection{Transactional Memory}

\begin{remark}
Synchronization with locks is hard to implement, not composable and pessimistic. The idea behind TM is to remove the burden from the programmer and let the system enforce synchronization.
\end{remark}

\begin{definition}
In transactional memory the programmer declares sections of code as \textbf{atomic} (a \textbf{transaction}), and the system ensures this code block takes effect atomically.
\end{definition}

\begin{lemma}
Transactions run isolated, consistent and atomic: other threads can only observe the start or end state. \textbf{HTM} offers fast execution at limited size, \textbf{STM} offers larger transactions at higher cost.
\end{lemma}

\begin{definition}
For each transaction: the \textbf{read set} is the set of variables read, the \textbf{write set} is the set of variables written.
\end{definition}

\begin{remark}
Conflicts occur when the write set of transaction $A$ overlaps with the read or write set of transaction $B$.
\end{remark}

\begin{remark}
TM can be implemented in hardware (\textbf{HTM}, using cache coherence infrastructure) or software (\textbf{STM}).
\end{remark}

\begin{definition}
\textbf{Weak isolation} only isolates transactions from each other.
\textbf{Strong isolation} isolates transactions from each other and from operations outside the transaction.
\end{definition}

\begin{definition}
\textbf{Flat Nesting}: inner transactions are invisible to the outside; if an inner transaction aborts, the outer transaction aborts too.
\end{definition}

\subsection{Distributed Memory and Message Passing}

\begin{remark}
Shared mutable data is difficult to handle. Alternative: local mutable data where threads send data to each other via \textbf{message passing}.
\end{remark}

\begin{definition}
\textbf{Message Passing} can be:
\begin{itemize}
    \item \textbf{Synchronous}: Sender blocks until the message is received.
    \item \textbf{Asynchronous}: Fire and forget (messages are buffered).
\end{itemize}
\end{definition}

\begin{definition}
\textbf{MPI (Message Passing Interface)}: Standard for distributed memory programming. Processes are grouped into \textbf{communicators}, each with a unique \textbf{process rank}.
\end{definition}

\begin{definition}
Message passing operations can be:
\begin{itemize}
    \item \textbf{Blocking}: Operation blocks until data is safely sent/received.
    \item \textbf{Non-blocking}: Returns immediately, allowing overlap between computation and communication.
\end{itemize}
\end{definition}

\begin{definition}
MPI collective operations:
\begin{itemize}
    \item \textbf{Barrier}: Synchronize all processes
    \item \textbf{Broadcast}: Send data from one process to all
    \item \textbf{Reduce}: Perform operation on data from all processes, result to one process
    \item \textbf{Scatter}: Distribute data in chunks from one process to all
    \item \textbf{Gather}: Collect data chunks from all processes into one
\end{itemize}
\end{definition}

\subsubsection{Actor Model}

\begin{definition}
\textbf{Actor Model}: A model of distributed computing where actors are the fundamental units of computation. Actors communicate by sending messages.
\end{definition}

\begin{remark}
Actors are isolated from each other and can only communicate by sending messages. This eliminates the need for locks or other synchronization primitives.
\end{remark}
